% GENERATED FILE. Edit canonical lessons, then run npm run generate:books.
% canonical-source-hash: fnv1a64:99b5b02e82913984
% canonical-lessons: GU-C373-upay, GU-C373-nirakaran, GU-C373-vikalp, GU-C373-muskeli, GU-C373-nadvu

\chapter{A Way Out}
\label{ch:gu-a-way-out}
\hlchaptermodality{373}

\begin{chapteropening}
\textbf{By the end of this chapter:} \emph{I can suggest what to do if something gets in the way.}

Complete the last lesson of chapter 373: I can suggest what to do if something gets in the way.
\end{chapteropening}

\section[\texorpdfstring{\gu{ઉપાય} (upāy)}{ઉપાય (upāy)}]{\gu{ઉપાય} (upāy) — a way out}

\label{lesson:GU-C373-upay}



\begin{quote}
Before the new one: say the Gujarati for alert, then the Gujarati for to include.
\end{quote}

\subsection*{You'll want to know: \gu{ઉપાય}}
\textbf{\gu{ઉપાય}} — \emph{upāy} — \textquotedblleft{}a way out\textquotedblright{}.

\begin{quote}
\gu{જો} \gu{આ} \gu{ઉપાય} \gu{કામ} \gu{કરશે} \gu{તો} \gu{તાવ} \gu{જલદી} \gu{ઊતરશે}.
\end{quote}

\begin{quote}
\emph{jo ā upāy kām karśe to tāv jaldī ūtarśe.}
\end{quote}

\begin{quote}
If this remedy works, the fever will soon come down.
\end{quote}

A masculine noun.

\begin{cousinweb}
\textbf{\gu{ઉપાય}} (\emph{upāy}) is a tatsama, borrowed from Sanskrit \emph{upāya}, \textquotedblleft{}coming near, a means, a way\textquotedblright{} (Monier-Williams).
\end{cousinweb}

\begin{grammarlens}[title={What you've built}]
One more word for saying what happens if.
\end{grammarlens}

\subsection*{Your turn}
\begin{itemize}
\raggedright
  \item \emph{Say it:} \emph{upāy}
  \item \emph{Say it:} \emph{upāy}, once more
  \item \emph{Say it:} \emph{jo ā upāy kām karśe to tāv jaldī ūtarśe.} — If this remedy works, the fever will soon come down.
  \item \emph{Recall:} say the Gujarati for alert, then the Gujarati for to include, then say \emph{upāy} again
\end{itemize}

\begin{tcolorbox}[breakable,colback=teal!4,colframe=teal!35!black,title={Before you move on}]
What does \gu{ઉપાય} mean? (\textquotedblleft{}A way out\textquotedblright{}.) Say it once more.
\end{tcolorbox}
\section[\texorpdfstring{\gu{નિરાકરણ} (nirākaraṇ)}{નિરાકરણ (nirākaraṇ)}]{\gu{નિરાકરણ} (nirākaraṇ) — a solution (to a problem); a settling}

\label{lesson:GU-C373-nirakaran}



\begin{quote}
Before the new one: say the Gujarati for to include, then the Gujarati for a way out.
\end{quote}

\subsection*{You'll want to know: \gu{નિરાકરણ}}
\textbf{\gu{નિરાકરણ}} — \emph{nirākaraṇ} — \textquotedblleft{}a solution (to a problem); a settling\textquotedblright{}.

\begin{quote}
\gu{જો} \gu{બંને} \gu{બાજુ} \gu{વાત} \gu{કરશે} \gu{તો} \gu{આ} \gu{વિવાદનું} \gu{નિરાકરણ} \gu{થશે}.
\end{quote}

\begin{quote}
\emph{jo banne bāju vāt karśe to ā vivādnũ nirākaraṇ thaśe.}
\end{quote}

\begin{quote}
If both sides talk, this dispute will be settled.
\end{quote}

A neuter noun.

\begin{cousinweb}
\textbf{\gu{નિરાકરણ}} (\emph{nirākaraṇ}) is a tatsama, borrowed from Sanskrit \emph{nirākaraṇa}, \textquotedblleft{}driving away, removing\textquotedblright{} (Monier-Williams). Gujarati uses it for clearing a problem away: its solution, a sense the source does not give.
\end{cousinweb}

\begin{grammarlens}[title={What you've built}]
One more word for saying what happens if.
\end{grammarlens}

\subsection*{Your turn}
\begin{itemize}
\raggedright
  \item \emph{Say it:} \emph{nirākaraṇ}
  \item \emph{Say it:} \emph{nirākaraṇ}, once more
  \item \emph{Say it:} \emph{jo banne bāju vāt karśe to ā vivādnũ nirākaraṇ thaśe.} — If both sides talk, this dispute will be settled.
  \item \emph{Recall:} say the Gujarati for to include, then the Gujarati for a way out, then say \emph{nirākaraṇ} again
\end{itemize}

\begin{tcolorbox}[breakable,colback=teal!4,colframe=teal!35!black,title={Before you move on}]
What does \gu{નિરાકરણ} mean? (\textquotedblleft{}A solution (to a problem); a settling\textquotedblright{}.) Say it once more.
\end{tcolorbox}
\section[\texorpdfstring{\gu{વિકલ્પ} (vikalp)}{વિકલ્પ (vikalp)}]{\gu{વિકલ્પ} (vikalp) — an alternative}

\label{lesson:GU-C373-vikalp}



\begin{quote}
Before the new one: say the Gujarati for a way out, then the Gujarati for a solution.
\end{quote}

\subsection*{You'll want to know: \gu{વિકલ્પ}}
\textbf{\gu{વિકલ્પ}} — \emph{vikalp} — \textquotedblleft{}an alternative\textquotedblright{}.

\begin{quote}
\gu{જો} \gu{ટ્રેન} \gu{રદ} \gu{થશે} \gu{તો} \gu{બસ} \gu{એક} \gu{વિકલ્પ} \gu{છે}.
\end{quote}

\begin{quote}
\emph{jo ṭren rad thaśe to bas ek vikalp che.}
\end{quote}

\begin{quote}
If the train is cancelled, the bus is an alternative.
\end{quote}

A masculine noun.

\begin{cousinweb}
\textbf{\gu{વિકલ્પ}} (\emph{vikalp}) is a tatsama, borrowed from Sanskrit \emph{vikalpa}, \textquotedblleft{}an alternative, an option\textquotedblright{} (Monier-Williams).
\end{cousinweb}

\begin{grammarlens}[title={What you've built}]
One more word for saying what happens if.
\end{grammarlens}

\subsection*{Your turn}
\begin{itemize}
\raggedright
  \item \emph{Say it:} \emph{vikalp}
  \item \emph{Say it:} \emph{vikalp}, once more
  \item \emph{Say it:} \emph{jo ṭren rad thaśe to bas ek vikalp che.} — If the train is cancelled, the bus is an alternative.
  \item \emph{Recall:} say the Gujarati for a way out, then the Gujarati for a solution, then say \emph{vikalp} again
\end{itemize}

\begin{tcolorbox}[breakable,colback=teal!4,colframe=teal!35!black,title={Before you move on}]
What does \gu{વિકલ્પ} mean? (\textquotedblleft{}An alternative\textquotedblright{}.) Say it once more.
\end{tcolorbox}
\section[\texorpdfstring{\gu{મુશ્કેલી} (muśkelī)}{મુશ્કેલી (muśkelī)}]{\gu{મુશ્કેલી} (muśkelī) — a difficulty}

\label{lesson:GU-C373-muskeli}



\begin{quote}
Before the new one: say the Gujarati for a solution, then the Gujarati for an alternative.
\end{quote}

\subsection*{You'll want to know: \gu{મુશ્કેલી}}
\textbf{\gu{મુશ્કેલી}} — \emph{muśkelī} — \textquotedblleft{}a difficulty\textquotedblright{}.

\begin{quote}
\gu{જો} \gu{કોઈ} \gu{મુશ્કેલી} \gu{આવશે} \gu{તો} \gu{મને} \gu{ફોન} \gu{કરો}.
\end{quote}

\begin{quote}
\emph{jo koī muśkelī āvśe to mane phon karo.}
\end{quote}

\begin{quote}
If any difficulty comes up, phone me.
\end{quote}

A feminine noun.

\begin{cousinweb}
\textbf{\gu{મુશ્કેલી}} (\emph{muśkelī}) is usually traced to Persian \emph{mushkil}, \textquotedblleft{}difficult\textquotedblright{}, a word of Arabic origin, with the Gujarati ending \emph{-ī}.
\end{cousinweb}

\begin{grammarlens}[title={What you've built}]
One more word for saying what happens if.
\end{grammarlens}

\subsection*{Your turn}
\begin{itemize}
\raggedright
  \item \emph{Say it:} \emph{muśkelī}
  \item \emph{Say it:} \emph{muśkelī}, once more
  \item \emph{Say it:} \emph{jo koī muśkelī āvśe to mane phon karo.} — If any difficulty comes up, phone me.
  \item \emph{Recall:} say the Gujarati for a solution, then the Gujarati for an alternative, then say \emph{muśkelī} again
\end{itemize}

\begin{tcolorbox}[breakable,colback=teal!4,colframe=teal!35!black,title={Before you move on}]
What does \gu{મુશ્કેલી} mean? (\textquotedblleft{}A difficulty\textquotedblright{}.) Say it once more.
\end{tcolorbox}
\section[\texorpdfstring{\gu{નડવું} (naḍvũ)}{નડવું (naḍvũ)}]{\gu{નડવું} (naḍvũ) — to get in the way}

\label{lesson:GU-C373-nadvu}



\begin{quote}
Before the new one: say the Gujarati for an alternative, then the Gujarati for a difficulty.
\end{quote}

\subsection*{You'll want to know: \gu{નડવું}}
\textbf{\gu{નડવું}} — \emph{naḍvũ} — \textquotedblleft{}to get in the way\textquotedblright{}.

\begin{quote}
\gu{જો} \gu{વરસાદ} \gu{નડશે} \gu{તો} \gu{મેચ} \gu{મુલતવી} \gu{રહેશે}.
\end{quote}

\begin{quote}
\emph{jo varsād naḍśe to mec mulatvī raheśe.}
\end{quote}

\begin{quote}
If the rain gets in the way, the match will be put off.
\end{quote}

Past (it …): \gu{તે} \gu{ભીડ} \gu{નડી}, \emph{te bhīḍ naḍī}.

\begin{cousinweb}
\textbf{\gu{નડવું}} (\emph{naḍvũ}) CDIAL lists the Gujarati verb, \textquotedblleft{}to obstruct\textquotedblright{}, under a reconstructed \emph{naṭati}, \textquotedblleft{}trembles, totters\textquotedblright{} (CDIAL 6934).
\end{cousinweb}

\begin{grammarlens}[title={What you've built}]
One more word for saying what happens if.
\end{grammarlens}

\subsection*{Your turn}
\begin{itemize}
\raggedright
  \item \emph{Say it:} \emph{naḍvũ}
  \item \emph{Say it:} \emph{naḍvũ}, once more
  \item \emph{Say it:} \emph{jo varsād naḍśe to mec mulatvī raheśe.} — If the rain gets in the way, the match will be put off.
  \item \emph{Recall:} say the Gujarati for an alternative, then the Gujarati for a difficulty, then say \emph{naḍvũ} again
\end{itemize}

\begin{tcolorbox}[breakable,colback=teal!4,colframe=teal!35!black,title={Before you move on}]
What does \gu{નડવું} mean? (\textquotedblleft{}To get in the way\textquotedblright{}.) Say it once more.
\end{tcolorbox}
